\documentclass{article}
\usepackage{amsmath,amssymb}
\begin{document}
\title{JOINT-FRONTIER1: Synchronization and Resource Accounting}
\date{2026-10-10}\maketitle
\paragraph{Status.} P versus NP remains OPEN. The preceding local experiment explored finite Pareto frontiers for certified CNF reductions, but did not establish machine time-space lower bounds.
\paragraph{Exact declared model.} States are canonical explicit CNFs; transitions are unit propagation, subsumption and Davis--Putnam elimination. Artificial positive weights are $1+m$ for unit propagation, $1+m^2$ for subsumption, and $1+m+|P_x||N_x|$ for variable elimination. Path time is the sum of these weights, and peak space is maximum number of explicit clauses. These are not measured RAM bit operations or peak workspace bits.
\paragraph{Recorded preceding computation.} The local run reported 32 CNFs, 16420 certified transitions, 2997 nondominated labels and 0.474 seconds. A 500-state exploration cap means these are frontiers of explored subgraphs, not necessarily global frontiers.
\paragraph{Correctness scope.} SAT-preserving transitions can be checked exhaustively on small instances. For a fixed finite graph, Pareto label pruning is exact only when all future transitions depend solely on the current state and the accumulated time and peak resources, with nonnegative costs. This does not transfer to unrestricted Boolean computation.
\paragraph{Next audit.} Verify each local artifact's exact content and hashes, rerun the code under a 30-second bound, and add encoded-bit and transient workspace accounting before claiming physical resource bounds. The all-input, source-first SAT successor remains OPEN.
\end{document}
